% Folder 91, batch 05 (archivists' pages 81-100).
% Pass: Opus 5.5 (claude-opus-5-5), 2026-10-03 — first pass, unchecked against the pages by a human.
%
% Even pages carry the notes; odd pages are typescript versos or blank and are skipped.
\documentclass[11pt,a4paper]{article}
\input{../preamble/grothendieck}

\folder{91}
\batch{5}
\pages{81}{100}
\foldertitle{Autour de Néron / Greenberg-Néron. Foncteurs Hom (méthodes non-projectives) : notes manuscrites (s.d.), lettre (1967).}
\dating{[à partir de 1964-vers 1970]}
\watermark{Édition de démonstration}

\begin{document}

\section{Parties constructibles, applications constructibles, ensembles élémentaires}

\page{81}
\note{le feuillet commence au milieu d'un argument sur les parties constructibles, venu d'avant le lot}
$Z$ constructible dans $X$ \uncertain{par définition} \add{\uncertain{donc}}
\uncertain{et} \uncertain{loc.} \ill{} \uncertain{localement} \uncertain{constructible} dans~$X$.

\struck{$Z$ loc. constructible}
\[
\begin{cases}
Z \text{ quasi-compact et} \\
\forall x \in Z,\ \exists \text{ un voisinage } U_x \text{ de } x \text{ dans } X,
\end{cases}
\]
\uncertain{tel que} $Z \cap U_x$ \uncertain{constr.} dans $U_x$
[\uncertain{donc} $\forall\, U' \subset U$, $U' \ni x$, \ill{} $= \text{\uncertain{ouvert}}$].

Réciproque évidente [\uncertain{car} $Z$ \uncertain{est} \uncertain{constructible} \uncertain{dans} $X$, \uncertain{si} \ill{} \uncertain{pour} \ill{} $U = X$].

\uncertain{Démonstration}. $Z$ est recouvert par \ill{} \ill{} \uncertain{finis} \uncertain{des}
$U_x$, \uncertain{soit} $U_i$. \uncertain{Les} $U_i \cap Z$ \uncertain{sont} \uncertain{constructibles}
dans $U_i$, \uncertain{donc} dans $X$, \uncertain{et} $Z = \bigcup Z_i$ \uncertain{est} \uncertain{constr.}

\textbf{Corollaire}. \uncertain{Supposons} $Z$ \uncertain{loc.} \ill{}, \ill{} \ill{}
\uncertain{tel que} \ill{}, $Z$ \uncertain{est} \uncertain{constructible}. \ill{}
\ill{} \ill{} \uncertain{et} \ill{} $x \in Z$. \uncertain{Il} \uncertain{existe} un
\uncertain{voisinage} \ill{} $U$ \uncertain{de} $x$ \uncertain{dans} $X$ \uncertain{tel que} $Z \cap U$ \uncertain{soit} \struck{\uncertain{fermé}} \add{\uncertain{fermé}}
\uncertain{et} \uncertain{constructible} dans~$U$,
[i.e. $U - Z \cap U$ \uncertain{quasi-compact} \add{(\uncertain{pour} $U$ \uncertain{quasi-compact}, \uncertain{et} \ill{})}
\ill{} \uncertain{l'intersection} \uncertain{de} \ill{} \uncertain{quasi-compacts}, \ill{} \ill{} \uncertain{dans}
\ill{} \uncertain{ouverts} \ill{} \ill{}\,].

\textbf{Exemple}. \uncertain{Si} $X$ \uncertain{est} \uncertain{un} \uncertain{préschéma} \uncertain{quasi-séparé} \uncertain{quasi-compact}
(\uncertain{p.\ ex.} un \uncertain{schéma} \uncertain{quasi-compact}, \uncertain{ou} un \uncertain{schéma} \uncertain{affine}) \struck{\ill{}}
\uncertain{alors}~: $Z$ constructible $\Longleftrightarrow$ $Z$ \uncertain{quasi-compact}, et
$Z$ loc. fermé dans $X$, \uncertain{donc} $\forall x \in Z$ \uncertain{un} \uncertain{voisinage} \ill{} \uncertain{affine} $U$ \uncertain{tel que} $Z \cap U$
\uncertain{soit} \uncertain{défini} \uncertain{par} \uncertain{un} \uncertain{nombre} \uncertain{fini} \uncertain{d'équations}.

\page{83}
\drawing{en haut à gauche, un croquis~: une grande région marquée « \uncertain{qu.~cpt} », contenant deux ouverts qui se recoupent, notés $U_f$ et $U_g$, avec leur intersection notée \uncertain{$U_f \cap U_g$}~; un petit cercle au-dessus, et un cercle barré plus à droite~; un $\overline{X}$ isolé.}

[\uncertain{Supposons} \uncertain{que} \uncertain{tout} \uncertain{pt} \uncertain{de} $X$ \uncertain{a} un \uncertain{syst.}
\uncertain{fond.} \uncertain{de} \uncertain{voisinages} \struck{\uncertain{ouverts}}\add{\uncertain{ouverts}}, \uncertain{quasi-compacts} \uncertain{dans} $X$.
[\uncertain{Alors} \ill{} \ill{} \ill{} \uncertain{ouvert} \uncertain{de} $X$ \uncertain{est} \uncertain{quasi-compact}]]
\note{passage encadré, en partie barré~; lecture très incertaine}

\marginal{$Z \subset U$ \uncertain{(schéma vertical)}}
Soit $T$ une partie \uncertain{fermée} \uncertain{irréductible} de~$X$.
\uncertain{Conditions} \uncertain{équivalentes}~:
\begin{enumerate}
  \item[(i)] Il existe une partie $Z$, \uncertain{partout} \uncertain{dense} \uncertain{dans} $T$,
  telle que $Z$ soit \uncertain{constructible} \uncertain{dans}~$X$~;
  \item[(ii)] Il existe un ouvert \add{\uncertain{quasi-compact}} $U$ de $X$, tel que
  $U \cap T \neq \emptyset$, et $U - U \cap T$ \ill{}
\end{enumerate}
\note{le feuillet, déchiré en bas, s'arrête sur cette ligne}

\page{86}
Préschémas atomiques~: (\uncertain{sommes} \uncertain{d'un} \uncertain{nombre} \uncertain{fini} \uncertain{de} \uncertain{spectres} de corps)
[\uncertain{sous-catégorie} \add{\uncertain{pleine} \uncertain{de}} (\uncertain{Sch}),]

\begin{tikzcd}
\text{Préschémas atomiques} \arrow[d, "{\text{foncteur pleinement fidèle}}"] \\
\text{Foncteurs covariants}\ (\mathrm{Corps}) \to (\mathrm{Ens})
\end{tikzcd}

i.e. \uncertain{contravariants} $(\text{Sch ponctuels})^{\circ} \to (\mathrm{Ens})$\note{« Sch ponctuels » est une lecture incertaine}.

\begin{tikzcd}
\text{Préschémas atomiques sur } k\ [\text{i.e. sur } \operatorname{Spec}(k)] \arrow[d, "{\text{foncteur pleinement fidèle}}"] \\
\text{Foncteurs covariants}\ (\text{Extensions de } k) \to (\mathrm{Ens})
\end{tikzcd}

\note{dans ces deux diagrammes, la lecture des étiquettes « foncteur pleinement fidèle » est incertaine, ainsi que « covariants » et « Préschémas »~; après la première étiquette, un mot illisible}

\marginal{\struck{$(\mathrm{At})$}$/\operatorname{Spec} k$}

Foncteurs \uncertain{''} $(\text{Sch ponctuels}/\operatorname{Spec} k)^{\circ} \to (\mathrm{Ens})$.

Soit $X$ un préschéma, il \uncertain{définit} \uncertain{un} \uncertain{préschéma} \uncertain{atomique}
$|X| = \coprod_{x \in X} \operatorname{Spec} k(x)$, \uncertain{donc} \struck{$(x_1, \dots, x_n)$} \ill{} \uncertain{un} \uncertain{foncteur} \uncertain{contravariant} \uncertain{défini} \uncertain{par}
$X \colon (\mathrm{Sch})^{\circ} \to (\mathrm{Ens})$ \struck{\ill{}} \add{\uncertain{admissible}} \uncertain{par} \uncertain{changement} \ill{}.

Bien \uncertain{sûr}, on obtient un foncteur $X \mapsto |X|$, \uncertain{de} la \uncertain{catégorie} des préschémas \uncertain{vers} \ill{}
\struck{\ill{}}~: \uncertain{il} \uncertain{n'est} \uncertain{pas} \uncertain{pleinement}
fidèle.

Soient $X$, $Y$ deux préschémas, une \uncertain{application} \uncertain{constructible}\note{« -plication » est souligné, et « pseudo-morphisme » plus haut sur la ligne~: lecture incertaine du rapport entre les deux mots} \uncertain{ou} \uncertain{pseudo-morphisme} de $X$ dans $Y$, \uncertain{est}
\uncertain{un} \uncertain{morphisme} $f \colon |X| \to |Y|$ qui a la propriété \struck{\ill{}}
suivante~: Pour \uncertain{toute} partie \struck{\uncertain{quasi-compacte}} \add{\uncertain{ouverte} \uncertain{affine}} $U$ de $X$,
il existe une partition de $U$ en des $U_i$ \add{\uncertain{loc.\ fermés}} (constructibles),
\uncertain{telle} \uncertain{que} $f \mid |U_i|$ soit un \uncertain{morphisme} $U_{i,\mathrm{red}} \to Y$.

\page{88}
\struck{Transitivité. [N.B. il suffit de le vérifier pour un \uncertain{ouvert} \uncertain{affine},
\uncertain{puis} \uncertain{quasi-compact} \add{$U^{\alpha}$} recouvrant $X$, et \uncertain{pour} $\forall\, U^{\alpha}$
\uncertain{trouver} un \uncertain{recouvrement} \uncertain{de} $U^{\alpha}$ \uncertain{par} \uncertain{des} \uncertain{parties} \add{\uncertain{loc.\ fermées}} constructibles
$U^{\alpha}_i$, telles que $f \mid |U^{\alpha}_i|$ soit \ill{} \uncertain{morphique}…}
\note{début de page biffé par deux grands traits en croix~; le crochet n'est pas refermé}

\marginal{\struck{Propositions}}

\uncertain{On} \uncertain{appelle} \uncertain{ensemble} \struck{\uncertain{algébrique}} \add{\uncertain{scindé}} \uncertain{élémentaire} un
préschéma \uncertain{quasi-séparé} \uncertain{quasi-compact} (\uncertain{quasi-noethérien} \uncertain{quasi-compact}),
\struck{\ill{}}, \uncertain{considéré} \uncertain{comme} \uncertain{objet} \uncertain{d'une} \uncertain{catégorie}
\uncertain{dont} les \uncertain{morphismes} \uncertain{sont} les applications constructibles.

\begin{tikzcd}[column sep=small, nodes={font=\scriptsize}]
(\mathrm{Sch}) \arrow[r] & (\mathrm{Sch}) \arrow[r, "\mathrm{red}"', "\text{fidèle}"] & (\text{ens. } \mathrm{Sch}) \arrow[r, "\text{fidèle}"'] & (\text{Sch. atom.}) \arrow[r, "\text{pl.\ fid.}"'] & {\underline{\operatorname{Hom}}((\mathrm{Corps}), \mathrm{Ens})}
\end{tikzcd}

\note{dans le premier terme, un mot biffé illisible après « Sch »~; les étiquettes « fidèle » et « pl.\ fid. » sont des lectures incertaines~; sous la chaîne, des arcs et des étiquettes « fidèle » répétées, en partie biffées par un zigzag~: on ne transcrit que la ligne principale}

\marginal{N.B. \uncertain{Dans} les \uncertain{foncteurs} \uncertain{ci-dessus}, \ill{} \uncertain{commutent} \uncertain{aux} $\varprojlim$ \uncertain{finies}}

Les \uncertain{définitions} \uncertain{donnent} \uncertain{aux} \uncertain{notions} \uncertain{de} \uncertain{les}
\uncertain{ensembles} \uncertain{élémentaires}, \uncertain{qui} \uncertain{se} \uncertain{définissent} \uncertain{par} \uncertain{deux}
\uncertain{d'ensembles} \uncertain{élémentaires} (\uncertain{i.e.} \uncertain{isom.} \uncertain{constructibles}).
\uncertain{Notons} \uncertain{comme} \ill{} \ill{}

\[
\begin{array}{ll}
X \to Y & \\
X_i & Y_i \\
X'_j & \\
X_i \cap X'_j \to Y_{ij} &
\end{array}
\]

[\struck{\textbf{Lemme}. Soit $X$ \uncertain{préschéma} \uncertain{quasi-compact} \uncertain{quasi-séparé}.
(\uncertain{$Z_i$}) une famille \uncertain{disjointe} de parties
constructibles \uncertain{recouvrant} $X$. Alors il existe une
\uncertain{partition} de $X$ en des parties \uncertain{loc.\ fermées} constructibles \uncertain{affines},
\uncertain{plus} \uncertain{fine} \uncertain{que} \uncertain{une} \uncertain{partition} \uncertain{de} $X$ \uncertain{par} des \uncertain{ens.\ constructibles}.}]

\textbf{Proposition}. Supposons $X$ quasi-séparé et quasi-compact, $f \colon |X| \to |Y|$ \uncertain{donnée}.
Conditions équivalentes~:
\begin{enumerate}
  \item[(i)] $f$ \uncertain{est} \uncertain{constructible}~;
  \item[(ii)] $\exists$ une partition \add{\uncertain{finie}} de $X$ en \uncertain{des} parties loc.\ fermées constructibles \add{\uncertain{affines}} $Z_i$,
  telles que $f \mid |Z_i|$ soit \uncertain{morphique}~;
  \item[(iii)] \struck{Il existe une \uncertain{suite} croissante d'ouverts quasi-compacts $U_i$ \uncertain{couvrant}
  $X$, \uncertain{telle} \uncertain{que} $U_{i+1} - U_i$ \ill{} soit \uncertain{morphique}}
\end{enumerate}
\note{la condition (iii) est biffée d'un zigzag~; la page s'arrête là}

\page{90}
\textbf{Proposition}. Soit $f \colon X \to Y$ \struck{\uncertain{un morphisme}} \uncertain{une} \uncertain{application} \uncertain{constructible}
d'\uncertain{ensembles} \uncertain{élémentaires}. \uncertain{Pour} \uncertain{toute} \uncertain{partie} \uncertain{constructible}
$E$ de $Y$, $f^{-1}(E)$ \uncertain{est} \uncertain{constructible}.

\marginal{$X \to Y$, $U$, $X_i$}\note{petit schéma griffonné en marge, en partie biffé}

\uncertain{En} \uncertain{effet}, soit $X = \coprod_i X_i$, \uncertain{les} $X_i$ \uncertain{loc.} \uncertain{fermés} \uncertain{et} \uncertain{constructibles}
dans $X$, \uncertain{tels} \uncertain{que} $f_i = f \mid |X_i|$ \uncertain{soit} \uncertain{un} \uncertain{morphisme}. \uncertain{Donc}
$f^{-1}(E) = \bigcup_i f_i^{-1}(E)$, \uncertain{et} \uncertain{il} \uncertain{suffit} \uncertain{de} \uncertain{prouver} \uncertain{que}
\uncertain{chaque} $f_i^{-1}(E) = f^{-1}(E) \cap X_i$ \uncertain{est} \uncertain{constructible} \uncertain{dans} $X_i$ \struck{\ill{}},
\uncertain{ce} \uncertain{qui} \uncertain{résulte} \uncertain{de} \ill{}, \uncertain{appliqué} \uncertain{aux} \uncertain{morphismes} \add{\uncertain{$f_i$}} $X_i \to X$ \uncertain{dans}
\uncertain{fini} \uncertain{et} \uncertain{des} \ill{} \uncertain{des} \uncertain{préschémas}, \uncertain{il} \uncertain{suffit}, \uncertain{par} \uncertain{une} \uncertain{structure} \uncertain{induite}
\uncertain{constructible} \uncertain{sur} $X_i$ \ill{}]
\ill{} \ill{} \ill{} \ill{}
\uncertain{un} \uncertain{des} \uncertain{morphismes} \add{(\uncertain{de} \uncertain{préschémas})}~; \uncertain{comme} \uncertain{il} \uncertain{est} \uncertain{connu}
\ill{} \ill{} \ill{}.
\marginal{\uncertain{toujours} \uncertain{l'argument}}

\textbf{Corollaire}. \uncertain{Isomorphisme} \uncertain{de} \uncertain{la} \uncertain{notion} \uncertain{de} \uncertain{partie}
constructible d'un \uncertain{ensemble} \uncertain{élémentaire}.

\textbf{Proposition}. Soit $f \colon X \to Y$ une application constructible
d'\struck{un \uncertain{ens. élémentaire}} \add{\uncertain{préschémas} \uncertain{quasi-compacts} \uncertain{quasi-séparés}}
\uncertain{d'un} \uncertain{ensemble} \uncertain{élémentaire}~; il \uncertain{faut} et il suffit qu'il existe
\uncertain{pour} un \uncertain{ensemble} \uncertain{fini} $I$, \uncertain{deux} \uncertain{familles} $(X_i)_{i \in I}$ de $X$, $(Y_i)_{i \in I}$ \uncertain{de} $Y$ \uncertain{de} \uncertain{parties} \uncertain{loc.} \uncertain{fermées}
\uncertain{constructibles} [\uncertain{formant} \uncertain{des} \uncertain{partitions}], \uncertain{telles} \uncertain{que} $f$ \uncertain{induise} \uncertain{un} \uncertain{isomorphisme}
de $X_i$ \uncertain{sur} $Y_i$.

\begin{tikzcd}
X \arrow[r, bend left=20, "f"] & Y \arrow[l, bend left=20, "g"]
\end{tikzcd}

\uncertain{C'est} \uncertain{évident}. \uncertain{Soit} $g$ l'\uncertain{inverse} \uncertain{de} $f$, \ill{} \ill{}, \ill{} \uncertain{la} \uncertain{nécessité}
\uncertain{dans} \uncertain{les} \ill{} \uncertain{de} $f$, $X_\alpha$, $Y_\beta$ \uncertain{des} \uncertain{partitions} \add{\uncertain{de} $X$, $Y$} \ill{}
\uncertain{de} \uncertain{conditions} \uncertain{constructibles} \uncertain{de} $f$ \uncertain{et} $g$.
\uncertain{Evidemment} \uncertain{dès} \uncertain{lors} \uncertain{les} \uncertain{hypothèses} \uncertain{donc} \ill{} \ill{}.
\note{la suite se lit au haut de la p.~92}

\page{92}
\note{suite de la démonstration de la proposition de la p.~90}
Soit $X_{\alpha\beta} = X_{\alpha} \cap f^{-1}(Y_{\beta})$,
$Y_{\beta\alpha} = Y_{\beta} \cap g^{-1}(X_{\alpha})$, \uncertain{alors}
$X_{\alpha\beta}$ et $Y_{\beta\alpha}$ \uncertain{sont} \uncertain{des} \uncertain{ensembles} \uncertain{constructibles}, \uncertain{qui} \uncertain{se}
\uncertain{correspondent} \uncertain{par} $f$, $g$.
\struck{\uncertain{Pour} \add{\ill{}} $X_{\alpha\beta}$ \uncertain{est} \ill{}}

\marginal{$X \leftarrow Y$, $Y_{\beta}$}

\struck{\uncertain{Comme} \ill{} \uncertain{fini} $\bigcup_{\lambda} X_{\alpha\beta\lambda}$, \uncertain{où} \uncertain{les} $X_{\alpha\beta\lambda}$ \uncertain{sont} \uncertain{des}
\uncertain{parties} \ill{} \uncertain{constructibles} \uncertain{de} \add{\ill{} $X_{\alpha\beta}$} $X_{\alpha\beta}$]. Alors
$g^{-1}(X_{\alpha\beta\lambda}) = g_{\beta}^{-1}(X_{\alpha\beta\lambda})$ \uncertain{est} \uncertain{un} \uncertain{ensemble} \ill{} \ill{}}
\uncertain{de} $Y_{\beta}$ \add{\uncertain{constructible}}.
\note{passage enfermé dans un trait sinueux qui semble le biffer~; lecture très incertaine}

\uncertain{D'ailleurs} $X_{\alpha\beta} = f_{\alpha}^{-1}(Y_{\beta})$ est \uncertain{un}
\struck{\uncertain{ensemble} \uncertain{élémentaire} \uncertain{de}} loc.\ fermé dans $X_{\alpha}$, donc dans $X$.

\uncertain{et} $f_{\alpha\beta} \colon X_{\alpha\beta} \to Y$ \struck{\uncertain{est}} \uncertain{un} \uncertain{morphisme} \add{\uncertain{de} \uncertain{schémas}} (\uncertain{car}
\uncertain{induit} \uncertain{par} $f_{\alpha} \colon X_{\alpha} \to Y$]. \struck{\ill{} \ill{}}
\uncertain{De} \uncertain{même} \uncertain{par} \uncertain{les} $Y_{\alpha\beta}$ \uncertain{sont}
\uncertain{constructibles}, \uncertain{et} \uncertain{loc.} fermés, \uncertain{et}
\add{$g_{\beta} \colon Y_{\alpha\beta} \to X$ \uncertain{est} \uncertain{un} \uncertain{morphisme} \uncertain{de} \uncertain{schémas}.}

\uncertain{Enfin} \uncertain{les} \uncertain{morphismes} $X_{\alpha\beta} \to Y_{\alpha\beta}$ \uncertain{et} \ill{} \uncertain{fixés} \uncertain{par} \uncertain{les}
\ill{}
$X_{\alpha\beta} \to Y_{\alpha\beta}$ [\uncertain{car} \uncertain{il} \uncertain{les} \uncertain{fixe} \uncertain{ainsi}], et
\uncertain{de} \uncertain{même} \uncertain{un} \uncertain{morphisme} $Y_{\alpha\beta} \to X_{\alpha\beta}$.
\uncertain{Ils} \uncertain{sont} \uncertain{inverses} \uncertain{l'un} \uncertain{de} \uncertain{l'autre}.

\textbf{Cor.~0}. Si $X$ \uncertain{noethérien}, \uncertain{alors} \ill{} $Y$ \uncertain{l'est}.

\textbf{Cor.~1}. \uncertain{Tout} un \uncertain{ensemble} \uncertain{élémentaire} \uncertain{est} \uncertain{isomorphe} \uncertain{à} \ill{} \uncertain{un}
$|X|$, \uncertain{avec} $X$ \uncertain{affine}…
\note{les Cor.~0 et Cor.~1 sont reliés par un grand trait courbe à la note de la marge inférieure}

\textbf{Cor.~2}. \uncertain{Tout} \uncertain{morphisme} \uncertain{d'ensembles} \uncertain{élémentaires}
\uncertain{est} \uncertain{isom.} \uncertain{à} un \uncertain{morphisme} \uncertain{d'ensembles} \uncertain{élémentaires}
\ill{} (\ill{} \uncertain{comme} \uncertain{application}
\uncertain{constructible}).

\struck{\uncertain{Dimension} \uncertain{d'un} \uncertain{ensemble} \uncertain{élémentaire}, \ill{} \uncertain{et} \ill{}.
\uncertain{Un} \uncertain{morphisme} $f \colon X \to Y$ \uncertain{de} \uncertain{préschémas} \uncertain{quasi-compacts} \uncertain{et} \uncertain{quasi-séparés}
\uncertain{d'un} \uncertain{ensemble} \uncertain{élémentaire} \uncertain{est} \uncertain{de type fini}…}
\note{ces lignes sont biffées de traits obliques}

\marginal{\uncertain{Notion} \uncertain{d'ens.} \uncertain{élémentaire} \uncertain{noethérien}~; [\uncertain{définition}~: \uncertain{$X$} \uncertain{noethérien}] \ill{} \textbf{Cor.~3}~: \ill{} \ill{}}
\note{note écrite en oblique dans la marge inférieure gauche, partiellement perdue (trous dans le papier)}

\page{94}
\struck{[\uncertain{quasi-fini}]}

\textbf{Morphismes de type fini} (\uncertain{d'ensembles} \uncertain{élémentaires})~: \uncertain{tel} \uncertain{qu'il}
\uncertain{existe} une partition de $X$ en \uncertain{ensembles} loc.\ fermés
\uncertain{constructibles} \uncertain{affines} $X_i$, telles que les
$f \mid X_i$ \uncertain{soient} \uncertain{des} \uncertain{morphismes} de type fini \struck{[\uncertain{quasi-finis}]}.
\marginal{\uncertain{et} \uncertain{c'est} \uncertain{indépendant} \uncertain{du} \uncertain{choix} \uncertain{des} \uncertain{préschémas} \uncertain{définis}~?}

\uncertain{Un} \uncertain{tel} \uncertain{morphisme} \uncertain{d'ens.} \uncertain{élémentaires} \uncertain{est} \uncertain{isom.}
\uncertain{à} \uncertain{un} \uncertain{morphisme} \uncertain{de} \uncertain{type} \uncertain{fini} \uncertain{de} \uncertain{schémas}
\uncertain{affines}, \uncertain{considéré} \uncertain{comme} \uncertain{morphisme} \uncertain{d'ens.}
\uncertain{élémentaires}). \textbf{Transitivité}.

\uncertain{Ens.} \uncertain{élémentaire} \uncertain{noethérien}, \uncertain{de} \uncertain{type} \uncertain{fini} \uncertain{sur}
un autre. \uncertain{Tous} \uncertain{sont} \uncertain{sur} $\mathbb{Z}$~!

\uncertain{Ens.} \uncertain{élémentaires} \uncertain{définis} \uncertain{par} une \uncertain{partie}
constructible d'un \uncertain{ens.} \uncertain{élémentaire}, \ill{}
\uncertain{immersion} \uncertain{canonique}, \ill{} \uncertain{de} $X$ \uncertain{en}
\uncertain{somme} $X = E \amalg (X - E)$, \ill{}
\uncertain{factorisabilité} \uncertain{par} une partie constructible, \uncertain{évidemment}~!

\uncertain{Produits} \uncertain{fibrés} [\uncertain{et} $\varinjlim$ finies] \uncertain{dans} \uncertain{la} \uncertain{catégorie}
des \uncertain{ens.} \uncertain{élémentaires}. \textbf{Existent}, et le
\struck{\uncertain{foncteur}} \struck{\uncertain{préschémas}} (\uncertain{préschémas} \uncertain{quasi-compacts} \uncertain{quasi-séparés}) $\to$
(\uncertain{ens.} \uncertain{élém.}) \uncertain{commute} \uncertain{aux} \uncertain{produits} \uncertain{fibrés}.
Si $X$ \uncertain{et} $Y$ \uncertain{sont} \uncertain{sur} $Z$, \uncertain{si} $X'$, $Y'$ \uncertain{sont} \uncertain{des}
parties constructibles \uncertain{de} $X$, $Y$, alors
\[
X' \times_Z Y' = \mathrm{pr}_1^{-1}(X') \cap \mathrm{pr}_2^{-1}(Y') \ \dots
\]

\page{96}
$(X_i)_{i \in I}$ \uncertain{système} \uncertain{projectif} \uncertain{filtrant} \uncertain{de}
\uncertain{préschémas} \struck{\uncertain{quasi-compacts}} \uncertain{quasi-compacts} \uncertain{quasi-séparés},
\uncertain{les} $f_{ij}$ \uncertain{étant} \uncertain{affines}. \uncertain{Je} \uncertain{suppose} \ill{} \ill{} \uncertain{finies}.

\uncertain{Donc} $X = \varprojlim X_i$ \uncertain{existe}.

\uncertain{A-t-on} \struck{\uncertain{élémentaires}}
\[
\operatorname{Cons} X \xrightarrow{\ \sim\ } \varprojlim \operatorname{Cons}(X_i) \quad ???
\]

\uncertain{Soit} $T$ un préschéma quasi-compact \uncertain{et} quasi-séparé,
\uncertain{alors} \ill{} \ill{}~:
\[
\operatorname{Hom}(\operatorname{Cons} T, \operatorname{Cons} X) = \operatorname{App\,cons}(T, X)
= \varinjlim_{\mathfrak{P}} \operatorname{Hom}\Bigl(\text{\struck{$\Pi$}}\ \coprod_{\alpha \in \operatorname{Ind}(\mathfrak{P})} T_{\alpha}^{\mathfrak{P}}, X\Bigr)
\]
\marginal{$\mathfrak{P}$ \uncertain{une} \uncertain{partition} \uncertain{finie} \uncertain{de} $T$ \uncertain{en} \uncertain{ens.} \uncertain{loc.} \uncertain{fermés} \uncertain{constructibles}}
\[
= \varinjlim_{\mathfrak{P}} \varprojlim_{i} \prod_{\alpha \in \operatorname{Ind}(\mathfrak{P})} \operatorname{Hom}(T_{\alpha}^{\mathfrak{P}}, X_i)
\]
\[
\operatorname{Hom}\bigl(\operatorname{Cons}(T), \text{\struck{$\operatorname{Cons} X$}}\ \varprojlim \operatorname{Cons}(X_i)\bigr)
= \varprojlim_{i} \varinjlim_{\mathfrak{P}} \prod_{\alpha \in \operatorname{Ind}(\mathfrak{P})} \operatorname{Hom}(T_{\alpha}^{\mathfrak{P}}, X_i)
\]
\note{le $\mathfrak{P}$ est un P gothique~; l'indice $\operatorname{Ind}(\mathfrak{P})$ est une lecture incertaine}

Cependant, considérons \uncertain{un} \uncertain{autre} \uncertain{système} \ill{}
$X_j$ \ill{} $(Y_j)_{j \in J}$, \uncertain{d'où}
$Y = \varprojlim Y_j$.
\uncertain{Quelles} \uncertain{sont} \uncertain{les} \struck{\uncertain{morphismes}} \uncertain{appl.} \uncertain{constructibles} \uncertain{de}
$X$ \uncertain{dans} $Y$~? \uncertain{Il} \uncertain{faut} \ill{} \add{\ill{}}
\[
\operatorname{Hom}(\operatorname{Cons} X, \operatorname{Cons} Y) = \varinjlim_{\mathfrak{P}} \prod_{\alpha} \operatorname{Hom}(\text{\struck{\ill{}}}\, X_{\alpha}^{\mathfrak{P}}, Y)
\]
$\mathfrak{P}$ \uncertain{partitions} \uncertain{finies} \uncertain{de} $X$ \uncertain{en} \uncertain{parties} \uncertain{loc.} \uncertain{fermées} \uncertain{constructibles} $X_{\alpha}^{\mathfrak{P}}$.
\note{la page s'arrête sur cette formule, sans conclure}

\page{98}
\note{suite de la formule qui clôt la p.~96}
\[
= \varinjlim_{\mathfrak{P}} \varprojlim_{j} \text{\struck{$\operatorname{Hom}$}}\ \prod_{\alpha} \operatorname{Hom}(X_{\alpha}^{\mathfrak{P}}, Y_j)
\]

\textbf{Ensembles quasi-élémentaires}.
\uncertain{Les} \uncertain{ens.} \uncertain{quasi-élémentaires} \uncertain{de} \uncertain{type} \uncertain{fini} \uncertain{sur} \uncertain{un}
\struck{\ill{}} \uncertain{ens.} \uncertain{quasi-élémentaire} \uncertain{noethérien}
\uncertain{forment} \uncertain{une} \uncertain{catégorie} \uncertain{géométrique} (?)

\page{100}
\struck{\uncertain{Pseudo-schémas} \uncertain{sur} $k$}

\begin{tikzcd}[column sep=large]
(\text{Ens. élémentaires sur } k) \arrow[r, "\text{pl.\ fidèle}"'] & {\underline{\operatorname{Hom}}\bigl((\text{Ens. élém.}/k)^{\circ}, (\mathrm{Ens})\bigr)}
\end{tikzcd}

\note{lecture incertaine de « Ens.\ élémentaires » et de « pl.\ fidèle » dans ce diagramme~; le $^{\circ}$ n'est pas sûr}

\begin{tikzcd}
& X & \\
S \arrow[ur, dashed] & S' \arrow[l] \arrow[u] & S'' \arrow[l, bend left=12] \arrow[l, bend right=12]
\end{tikzcd}

\struck{\uncertain{application} \uncertain{surjective}} \quad \uncertain{morphisme} \uncertain{fid.\ plat}
\note{au-dessous du diagramme de descente, « application surjective » est biffé et « morphisme fid.\ plat » écrit à côté, souligné~; la flèche $S' \to S$ est l'épimorphisme en question}

\textbf{\uncertain{Tout} \uncertain{épimorphisme} \uncertain{est} \uncertain{un} \uncertain{morphisme} \uncertain{de} \uncertain{descente}}

\begin{enumerate}
  \item[1)] \uncertain{Existence} \uncertain{de} $\varprojlim$ \uncertain{finies}
  \item[2)] \uncertain{Existence} \uncertain{du} \uncertain{quotient} \uncertain{par} \uncertain{un} \uncertain{groupoïde} \uncertain{d'équivalence}
  \marginal{[\uncertain{dans} \uncertain{ce} \ill{} \uncertain{pour} \uncertain{groupoïdes} \uncertain{d'équivalence}…]}
  \item[3)] \uncertain{Tout} \uncertain{épimorphisme} \uncertain{est} \uncertain{effectif}
  \item[4)] \uncertain{Tout} \uncertain{épim.} \uncertain{est} \uncertain{universel}
  \item[\struck{5)}] \struck{\uncertain{Tout} \uncertain{morphisme} \uncertain{se} \uncertain{factorise} \uncertain{par} \uncertain{un} \uncertain{produit} \ill{} [\uncertain{épim.}, \ill{}]}
\end{enumerate}
\note{les points 1)--2) et 3)--4) sont groupés par des crochets, l'ensemble par une accolade~; le point 5) est biffé}

\begin{tikzcd}[column sep=small]
X \times_Y X \arrow[r, bend left=12] \arrow[r, bend right=12] & X \arrow[r] \arrow[dr] & Y \\
& & X/\mathrm{Rel}
\end{tikzcd}

\note{« $X/\mathrm{Rel}$ » est une lecture incertaine~; ce petit diagramme est entouré d'une boucle, reliée par un trait à la marge droite de la liste~; devant lui, une flèche biffée}

\begin{enumerate}
  \item[1')] \uncertain{Existence} \uncertain{de} $\varinjlim$ \uncertain{finies}
  \item[2')] \uncertain{Existence} \uncertain{du} \uncertain{conoyau} \uncertain{par} \uncertain{un} \uncertain{couple} \uncertain{d'équivalence}
  \item[3')] \uncertain{Tout} \uncertain{monom.} \uncertain{est} \uncertain{effectif}
  \item[4')] \uncertain{Changement} \uncertain{de} \uncertain{base} \uncertain{dans} \uncertain{un} \uncertain{monomorphisme}
\end{enumerate}

\textbf{Axiome géométrique}
\[
X \times \coprod_i Y_i \simeq \coprod_i X \times Y_i
\]
[\ill{}~: \uncertain{dire} \uncertain{que} \uncertain{le} \ill{} \uncertain{commute} \uncertain{aux} \ill{}
\uncertain{dans} \uncertain{les} $\varinjlim$ \uncertain{finies} ???]

\[
R \to S
\]
\note{la page s'achève sur ce « $R \to S$ » isolé~; la suite, si elle existe, est hors du lot}

\end{document}
